\chapter{Architectural Design}

\section{Procedure for approximating Error statistics on filter kernels}

One of the fundamental results in DSP is that the ideal window function
\[
w(t) = \mathrm{sinc}(t)
\]
in the time domain corresponds to a perfect rectangular pulse in the frequency domain~\cite{OppenheimDSP}. However, as the envelope $\tfrac{1}{x}$ decays the amplitude of $\sin(x)$ monotonically, the function has an infinite domain that is impossible to capture discretely. Thus, there is an associated truncation error $T(x)$, namely, it is the error between the tail \& partial power series. This implies the existence of an associated error term $\mathbb{E}(x)$ that depends on the truncation point $x$. The second contributor is \textit{coefficient quantization error}.

\subsection{Optimal Coefficient Quantization via Projection Methods}

Consider a function $f(x)$ represented as a polynomial series:
\begin{equation}
f(x) = \sum_{n=0}^{N} c_n x^{2n}
\end{equation}
where $c_n$ are the series coefficients. In fixed-point hardware implementations, these coefficients must be quantized to dyadic rationals of the form $\tilde{c}_n = a_n/2^{m_n}$, where $a_n \in \mathbb{Z}$ and $m_n$ determines the bit precision allocated to coefficient $n$.

The quantization error for coefficient $n$ is:
\begin{equation}
\epsilon_n = c_n - \tilde{c}_n
\end{equation}

The total approximation error introduced by coefficient quantization is:
\begin{equation}
E(x) = f(x) - \tilde{f}(x) = \sum_{n=0}^{N} \epsilon_n x^{2n}
\end{equation}

Our objective is to determine the optimal bit allocation vector $\mathbf{m} = [m_0, m_1, \ldots, m_N]$ that minimizes the projected error under a total bit budget constraint $\sum_{n=0}^{N} m_n \leq B_{\text{total}}$ for a LPF. Framing the problem in terms of an analytically continued series, $R(x)$, exposes a singular metric that packs in both 'phase' and 'bit' quantization error, expressed as a sensitivity matrix. This is effective on the computational front, allowing for a signal to be dually low-passed 
\& re-weighted according to tolerance error. As will be shown, it also a neccessery ingredient for autonomously reducing memory pressure via. a parameterized 'sacrifice on accuracy'.

\subsection{Projection Operator and Error Metric}

Rather than minimizing the error in the standard $L^2$ norm, we introduce a projection operator $R(x)$ that captures the application-specific error weighting. The projected error is defined as:
\begin{equation}
\mathcal{E}_{\text{proj}} = \langle E(x), R(x) \rangle = \int_{0}^{x_{\max}} E(x) \cdot \overline{R(x)} \, dx
\label{eq:projected_error}
\end{equation}
where $\overline{R(x)}$ denotes the complex conjugate of $R(x)$, and $\langle \cdot, \cdot \rangle$ represents the inner product in $L^2[0, x_{\max}]$.

Substituting the error polynomial:
\begin{equation}
\mathcal{E}_{\text{proj}} = \int_{0}^{x_{\max}} \left[\sum_{n=0}^{N} \epsilon_n x^{2n}\right] \overline{R(x)} \, dx = \sum_{n=0}^{N} \epsilon_n \int_{0}^{x_{\max}} x^{2n} \overline{R(x)} \, dx
\end{equation}

\subsection{Sensitivity Matrix Construction}

The sensitivity of the projected error to coefficient $n$ is defined as:
\begin{equation}
S_n = \left|\frac{\partial \mathcal{E}_{\text{proj}}}{\partial \epsilon_n}\right| = \left|\int_{0}^{x_{\max}} x^{2n} \overline{R(x)} \, dx\right|
\label{eq:sensitivity}
\end{equation}

This represents how strongly a unit error in coefficient $c_n$ affects the overall projected error. For a complex projector $R(x) = R_{\text{real}}(x) + iR_{\text{imag}}(x)$, we can decompose the sensitivity into real and imaginary components:
\begin{equation}
S_n^{\text{real}} = \left|\int_{0}^{x_{\max}} x^{2n} R_{\text{real}}(x) \, dx\right|, \quad
S_n^{\text{imag}} = \left|\int_{0}^{x_{\max}} x^{2n} R_{\text{imag}}(x) \, dx\right|
\end{equation}

The \textit{phase sensitivity} and \textit{magnitude sensitivity} can be defined as:
\begin{align}
S_n^{\text{phase}} &= S_n^{\text{imag}} \label{eq:phase_sens} \\
S_n^{\text{mag}} &= S_n^{\text{real}} \label{eq:mag_sens}
\end{align}

The complete sensitivity matrix is:
\begin{equation}
\mathbf{S} = \begin{bmatrix}
S_0^{\text{mag}} & S_0^{\text{phase}} \\
S_1^{\text{mag}} & S_1^{\text{phase}} \\
\vdots & \vdots \\
S_N^{\text{mag}} & S_N^{\text{phase}}
\end{bmatrix}
\end{equation}

\subsection{Quantization Error Bounds}

For dyadic quantization with $m_n$ bits allocated to coefficient $n$, the quantization error is bounded by:
\begin{equation}
|\epsilon_n| \leq \frac{1}{2^{m_n+1}}
\label{eq:quant_bound}
\end{equation}

The worst-case projected error is therefore:
\begin{equation}
|\mathcal{E}_{\text{proj}}| \leq \sum_{n=0}^{N} S_n |\epsilon_n| \leq \sum_{n=0}^{N} \frac{S_n}{2^{m_n+1}}
\label{eq:error_bound}
\end{equation}

\subsection{Optimal Bit Allocation Strategy}

\subsubsection{Method 1: Sensitivity-Proportional Allocation}

Allocate bits proportionally to sensitivity:
\begin{equation}
m_n = m_{\text{base}} + \left\lfloor \alpha \cdot \frac{S_n}{S_{\max}} \right\rfloor
\label{eq:prop_alloc}
\end{equation}
where $S_{\max} = \max_n S_n$, $m_{\text{base}}$ is the minimum bit allocation, and $\alpha$ controls the dynamic range of bit allocation (typically $\alpha \in [3, 8]$).

The allocation is subject to the constraint:
\begin{equation}
\sum_{n=0}^{N} m_n = B_{\text{total}}
\end{equation}

\subsubsection{Method 2: Weighted Phase-Magnitude Allocation}

For applications requiring different emphasis on phase versus magnitude accuracy, define a weighted sensitivity:
\begin{equation}
S_n^{\text{weighted}} = \sqrt{w_{\text{mag}} (S_n^{\text{mag}})^2 + w_{\text{phase}} (S_n^{\text{phase}})^2}
\label{eq:weighted_sens}
\end{equation}
where $w_{\text{mag}}$ and $w_{\text{phase}}$ are application-dependent weights satisfying $w_{\text{mag}} + w_{\text{phase}} = 1$.

\subsubsection{Method 3: Greedy Optimization}

For the global optimum under constraint \eqref{eq:error_bound}, we employ a greedy algorithm. The algorithm proceeds as follows:

\medskip
\noindent\textbf{Greedy Bit Allocation Algorithm:}
\begin{enumerate}
\item Initialize $m_n = 1$ for all $n \in \{0, \ldots, N\}$
\item Set $B_{\text{remaining}} = B_{\text{total}} - (N+1)$
\item While $B_{\text{remaining}} > 0$:
    \begin{enumerate}
    \item Compute the error reduction for adding one bit to coefficient $n$:
    \begin{equation}
    \Delta_n = S_n \left(\frac{1}{2^{m_n+1}} - \frac{1}{2^{m_n+2}}\right) \quad \text{for all } n
    \end{equation}
    \item Find the coefficient with maximum error reduction: $n^* = \arg\max_n \Delta_n$
    \item Increment the bit allocation: $m_{n^*} \gets m_{n^*} + 1$
    \item Decrement the remaining budget: $B_{\text{remaining}} \gets B_{\text{remaining}} - 1$
    \end{enumerate}
\item Return the optimal bit allocation vector $\mathbf{m} = [m_0, m_1, \ldots, m_N]$
\end{enumerate}

\medskip
\noindent This greedy approach is optimal for the additive error model in \eqref{eq:error_bound} and has computational complexity $O(B_{\text{total}} \cdot N)$.

\subsection{Evaluation of Projected Error}

After bit allocation, the actual projected error can be computed as:
\begin{equation}
\mathcal{E}_{\text{proj}} = \left|\sum_{n=0}^{N} (c_n - \tilde{c}_n) \int_{0}^{x_{\max}} x^{2n} \overline{R(x)} \, dx\right|
\end{equation}

For numerical evaluation, discretize the integral:
\begin{equation}
\mathcal{E}_{\text{proj}} \approx \left|\sum_{n=0}^{N} \epsilon_n \sum_{k=0}^{K-1} x_k^{2n} \overline{R(x_k)} \Delta x\right|
\end{equation}
where $x_k = k \cdot x_{\max} / K$ and $\Delta x = x_{\max} / K$.

The real and imaginary components of the error are:
\begin{align}
\mathcal{E}_{\text{real}} &= \text{Re}\left[\sum_{n=0}^{N} \epsilon_n \int_{0}^{x_{\max}} x^{2n} \overline{R(x)} \, dx\right] \\
\mathcal{E}_{\text{imag}} &= \text{Im}\left[\sum_{n=0}^{N} \epsilon_n \int_{0}^{x_{\max}} x^{2n} \overline{R(x)} \, dx\right]
\end{align}

\newpage
\subsection{Projector Derivation for Sinc Function}

The power series for $\mathrm{sinc}(x)$ is given as:
\[
\begin{aligned}
\mathrm{sinc}(x) &= \sum_{n=0}^\infty \frac{-1^n}{\Gamma(2n)}x^{2n + 1}
\end{aligned}
\]

But due to the necessity of considering truncated power series, introduce the upper bound $N$.Then re-indexing to $2N$ we obtain:
\[
\begin{aligned}
\mathrm{sinc}(x) &= \sum_{k = 0}^{2N} \frac{(-1)^{\left\lfloor k/2 \right\rfloor} x^{k+1}}{\Gamma(k)}
&= \sum_{k = 0}^{2N} \delta_{k \equiv 0 \, (\mathrm{mod}\,2)} \cdot \frac{i^k \cdot x^{k+1}}{\Gamma(k)}
\end{aligned}
\]

Implicitly allowing for the analytic continuation of $E$ ($E: \mathbb{R} \to \mathbb{R} \rightarrow E: \mathbb{C} \to \mathbb{C}$) due to the introduction of the \emph{Kronecker delta}\cite{}.This projector naturally selects even-indexed coefficients, which are precisely the non-zero terms in the sinc Taylor series. 

\medskip
Using the fact that we can exchange the sum \& integrand due to linearity\cite{} and using the contour representation of the kroenecker delta for term extraction\cite{} (I.e. as a \textit{cauchy kernel}) we derive:

\[
\begin{aligned}
\frac{1}{2\pi i}\oint_\gamma \sum_{k=0}^{2N} \frac{1}{\zeta(1 - \zeta)^2} \frac{i^k \cdot x^{k+1}}{\Gamma(k)} d\zeta
\end{aligned}
\]\eqref{}

Then let 
\[
\begin{aligned}
S &:= \sum_{k=0}^{2N} \frac{i^k \cdot x^{k+1}}{\Gamma(k)} \\
M &:= 2N \\
\implies S &= -i \sum_{j=1}^{M+1} \frac{(ix)^j}{\Gamma(j)} \\
&= -i \left( \exp(ix) - 1 - \sum_{j=M+2}^{\infty} \frac{(ix)^j}{\Gamma(j)} \right)
\end{aligned}
\]\eqref{}

We may then use the incomplete-gamma identity\cite{} for the power series of $e^x$ yielding:

\[
\begin{aligned}
S &= - i \left[ \exp(ix)\left(\frac{\Gamma(n+2, ix)}{\Gamma(n+2)} - 1\right) \right]
&= i - \frac{i\,\exp(ix)\,\Gamma(n+2, ix)}{\Gamma(n+2)}
\end{aligned}
\]\eqref{}

The projector identity*\cite{} finally gives:

\[
\begin{aligned}
\mathbb{P}(x) &= \frac{1}{2\pi i}\oint_\gamma \mathrm{B}(z)\mathrm{A}(\frac{x}{z})dz \\
\implies R(x) &= \frac{1}{2\pi \Gamma(M + 2)}\oint_\gamma \frac{\Gamma(M + 2) - \exp(\frac{ix}{z}) \Gamma(M + 2, \frac{ix}{z})}{z(1-z)^2}dz
\end{aligned}
\]\eqref{}

Which can be expressed more robustly using a change of variable $\frac{x}{z} \to \theta$

\[
\begin{aligned}
\frac{1}{2\pi \Gamma(K)}\oint_{C_\theta} (\Gamma(K) - e^{i\theta}\Gamma(K, i\theta))(\frac{\theta}{\theta^2 - x^2})d\theta
\end{aligned}
\]\eqref{}

Where C is the unit image circle.
\medskip

By the residue theorem, evaluating at the poles $\theta = \pm x$:
\[
\begin{aligned}
\frac{1}{2\pi \Gamma(K)}\oint_{C_\theta} \frac{\theta \Gamma(K)}{\theta^2 - x^2} d\theta = \frac{1}{2\Gamma(K)} \cdot 2\pi i \left[\text{Res}_{\theta=x} + \text{Res}_{\theta=-x}\right] \Gamma(K)
\end{aligned}
\]

The residues are:
\[
\begin{aligned}
\text{Res}_{\theta=x} \frac{\theta}{\theta^2 - x^2} &= \lim_{\theta \to x} \frac{\theta(\theta - x)}{(\theta-x)(\theta+x)} = \frac{x}{2x} = \frac{1}{2} \\
\text{Res}_{\theta=-x} \frac{\theta}{\theta^2 - x^2} &= \lim_{\theta \to -x} \frac{\theta(\theta + x)}{(\theta-x)(\theta+x)} = \frac{-x}{-2x} = \frac{1}{2}
\end{aligned}
\]

Therefore:
\begin{equation}
\frac{1}{2\pi \Gamma(K)}\oint_{C_\theta} \frac{\theta \Gamma(K)}{\theta^2 - x^2} d\theta = \frac{2\pi i}{2\Gamma(K)} \cdot \Gamma(K) \cdot 1 = i
\end{equation}

For the second integral involving the incomplete gamma function:
\begin{equation}
\frac{1}{2\pi \Gamma(K)}\oint_{C_\theta} \frac{\theta e^{i\theta}\Gamma(K, i\theta)}{\theta^2 - x^2} d\theta
\end{equation}

This can be evaluated by recognizing that for $\theta = x$:
\[
\begin{aligned}
\text{Res}_{\theta=x} \frac{\theta e^{i\theta}\Gamma(K, i\theta)}{\theta^2 - x^2} = \frac{x e^{ix}\Gamma(K, ix)}{2x} = \frac{e^{ix}\Gamma(K, ix)}{2}
\end{aligned}
\]

And for $\theta = -x$:
\[
\begin{aligned}
\text{Res}_{\theta=-x} \frac{\theta e^{i\theta}\Gamma(K, i\theta)}{\theta^2 - x^2} = \frac{-x e^{-ix}\Gamma(K, -ix)}{-2x} = \frac{e^{-ix}\Gamma(K, -ix)}{2}
\end{aligned}
\]

The exact solution to this part of the integral split is:
\[
\begin{aligned}
& i + \frac{1}{\Gamma(K)}\left[ \frac{ae^ix + be^ix}{2i} \right] \\
& a := \Gamma(k, ix) \\
& b := \Gamma(k, -ix)
\end{aligned}
\]\eqref{}
Where similarity to the complex exponential form of $cos(x)$ is noted.

This can be simplified using the properties of complex conjugation\eqref{} to:
\[
\begin{aligned}
i \left[ 1 - \frac{\Re(e^{ix}\Gamma(K, ix))}{\Gamma(K)} \right]
\end{aligned}
\]

\textit{Hartog's theorem}\cite{} further reduces the expression $\Re(e^{ix}\Gamma(K, ix))$ to:
\[
\begin{aligned}
\Re\left(\omega^{K}\sum_{j=0}^{\infty}\frac{\omega^{j}\,\Gamma(K)}{\Gamma(K+j+1)}\right) \\
&= \Re\left(\frac{\omega^K}{K}\sum_{j=1}^{\infty}\frac{\omega^{j}}{\Gamma(K+j)}\right) \\
&= \Re\left(\frac{\omega^K}{\Gamma(K+1)}\sum_{j=1}^{\infty}\frac{\omega^j}{K^{(j)}}\right)
\end{aligned}
\] Where $K^{(j)}$ is the rising factorial.
\medskip

Noticing that the outside term $\frac{\omega^K}{\Gamma(K+1)}$ will always be real-valued because $K \in \mathbb{N} \text{, } K \geq 1 \text{, } N \geq 1$ \& $K = M + 2 = 2N + 2 \implies K$ is even) means it can be moved outside. We subsequently conclude, using the linearity of $\Re$ over a series sum, that $\iff \text{j is even}$ will a term be emitted. Thus:

\[
\begin{aligned}
&= {\omega^K}{\Gamma(K+1)} \sum_{j=1}^{\infty} \frac{\omega^{2j}}{K^{(2j)}} \\
&= {\omega^K}{\Gamma(K+1)} \sum_{j=0}^{\infty} \frac{(-1)^{j+1}x^{2j+2}}{K^{(2j+2)}}
\end{aligned}
\]

This can be represented in terms of the hypogeometric series*\cite{}. Substituting values we finally arrive at the answer:

\[
\begin{aligned}
\boxed{i\!\left(1 - \frac{x^{2}\,\omega^{K}}{K+1}\,\mathbb{H}\right)} \\
\mathbb{H} := \pFq{1}{2}{1,\frac{K}{2}+1,\frac{K}{2}+\frac{3}{2}}{}{\frac{-x^2}{4}}
\end{aligned}
\]

Where for the special case $K=0 \rightarrow \mathbb{H} = \mathrm{sinc}(x)$.

\subsubsection{Approximation of the projector function}

We can forgo an exact solution for an approximate solution using the mirror property\cite{} of incomplete gamma functions \& assuming $x > 0 \text{, } x \in \mathbb{R}$. This is so that we can exploit the fact that the dominant contribution comes from the positive pole:

\begin{equation}
\frac{1}{2\pi \Gamma(K)}\oint_{C_\theta} \frac{\theta e^{i\theta}\Gamma(K, i\theta)}{\theta^2 - x^2} d\theta \approx \frac{2\pi i}{2\pi \Gamma(K)} \cdot \frac{e^{ix}\Gamma(K, ix)}{2} = \frac{i e^{ix}\Gamma(K, ix)}{2\Gamma(K)}
\end{equation}

Combining both terms from \eqref{}:
\[
\begin{aligned}
&\frac{1}{2\pi \Gamma(K)}\oint_{C_\theta} \left(\Gamma(K) - e^{i\theta}\Gamma(K, i\theta)\right)\left(\frac{\theta}{\theta^2 - x^2}\right)d\theta \nonumber \\
&= i - \frac{i e^{ix}\Gamma(K, ix)}{2\Gamma(K)}
\end{aligned}
\]\label{eq:final_unsimplified_result}

Simplifying:
\[
\begin{aligned}
&= i\left(1 - \frac{e^{ix}\Gamma(K, ix)}{2\Gamma(K)}\right)
\end{aligned}
\]

This can be rewritten in terms of the regularized incomplete gamma function $Q(K, ix) = \Gamma(K, ix)/\Gamma(K)$:
\begin{equation}
\boxed{
= i\left(1 - \frac{e^{ix}Q(K, ix)}{2}\right)
}\label{eq:final_result}
\end{equation}